\documentclass[11pt]{article}
\usepackage[margin=2.6cm]{geometry}
\usepackage{amsmath, amssymb, amsthm}
\usepackage[colorlinks=true, linkcolor=blue!60!black, citecolor=blue!60!black]{hyperref}
\usepackage{microtype}

\newtheorem{theorem}{Theorem}
\newtheorem{proposition}[theorem]{Proposition}
\newtheorem{lemma}[theorem]{Lemma}
\newtheorem{corollary}[theorem]{Corollary}
\newtheorem{conjecture}[theorem]{Conjecture}
\theoremstyle{definition}
\newtheorem{definition}[theorem]{Definition}
\newtheorem{remark}[theorem]{Remark}

\newcommand{\Gr}{\mathrm{Gr}}
\newcommand{\E}{\mathbb{E}}
\newcommand{\Var}{\mathrm{Var}}
\newcommand{\dG}{d_G}
\newcommand{\cX}{\mathcal{X}}
\newcommand{\dd}{\mathrm{d}}
\newcommand{\H}{\mathcal{H}}
\newcommand{\law}{\overset{d}{=}}
\newcommand{\xto}{\xrightarrow{d}}

\title{\textbf{Research Report: The Diversity Dichotomy as a Synchronization Transition}\\
\large Corrected Theory of Dependent U-Statistics in Grassmannian Swarm Dynamics\\
\small (Cycle 1 of the autonomous research program ``U-Dichotomie'')}
\author{Autonomous research agent (GPU-assisted: numerics + theory)}
\date{\today}

\begin{document}
\maketitle

\begin{abstract}
The formal problem statement conjectured that the observed variance dichotomy
$\Var(\Delta)=\mathcal O(N^{-2})$ vs.\ $\Var(\Delta)=\Theta(N^{-1})$ of the mean swarm
diversity is caused by a breaking of $O(d)$-invariance of the one-point marginal
$\pi_1$ once the consensus coupling $\eta>0$ is switched on. We show that this
mechanism is \emph{wrong}, for a structural reason: any $O(d)$-equivariant,
permutation-symmetric Markov dynamics with a unique stationary law has
$\pi_1=\mu_H$ (Haar) \emph{for every} $\eta\ge 0$; the symmetry breaking is
\emph{spontaneous} (a random cluster location restores the symmetry), and consequently
the Hoeffding projection $h_1\equiv0$ and the proposed effective variance
$\zeta_1^*\equiv 0$ identically. The dichotomy nevertheless exists and is real;
we prove that it is governed by the \emph{mean-field synchronization (bifurcation)
transition} at a critical coupling $\eta_c$: below $\eta_c$ the swarm is chaotic,
$\Var(\Delta)=2\zeta_2(\eta)/(N(N-1))=\Theta(N^{-2})$ and $N(\Delta-\theta)$ converges
to a weighted chi-square law; above $\eta_c$ the swarm forms a cluster of width
$\propto D/\eta$ and
$\Var(\Delta)=\sigma_\eta^2/N+\mathcal O(N^{-2})$ with
$\sigma_\eta^2=4\,c_1(\eta)+O(N^{-1})$, $c_1(\eta)=\E\bigl[\Var_{\nu_F}(m_F)\bigr]$
the variance of the \emph{conditional} first Hoeffding projection, so that
$\sqrt N(\Delta-\theta)\Rightarrow\mathcal N(0,\sigma_\eta^2)$ even though the classical
degeneracy $\zeta_1=0$ persists. A complete GPU Monte-Carlo campaign on
$\Gr(1,8)$ and $\Gr(2,8)$ confirms all exponents numerically
($-1.994/-2.013/-1.004/-1.001$), the Gaussian limit ($|$skew$|\lesssim0.13$,
kurtosis $\approx3$, QQ--$r\ge0.9988$), and the weighted chi-square limit in the Haar
case including the analytic constant $\sum_k\lambda_k^2=\pi^2/48$,
$\Var(N(\Delta-\theta))\to\pi^2/24$ on $\Gr(1,2)$. The susceptibility-like statistic
$N\,\Var(\Delta)$ jumps by a factor $\approx30$ at the transition, locating
$\eta_c\approx25$--$28$ for our discretization, and decays like $\propto\eta^{-1}$
in the ordered phase, as predicted by $c_1(\eta)\asymp D/(\kappa\eta)$.
\end{abstract}

% =====================================================================
\section{Setting}
Let $\cX=\Gr(k,d)$ with geodesic distance $\dG$ and let the swarm
$\mathbf V=(V_1,\dots,V_N)\in\cX^N$ evolve under the covariant consensus-SGD kernel
$P_\eta$ with consensus strength $\eta\ge0$ and isotropic noise of strength $D$.
We assume throughout:
\begin{enumerate}
\item[(A1)] \emph{Equivariance:} for every $R\in O(d)$,
$P_\eta(R\mathbf x,\cdot)=(\delta_{R})^{\times N} P_\eta(\mathbf x,R^{-1}\cdot)$.
This holds for the consensus dynamics because the consensus field
$\bar S=\frac1N\sum_j Q_jQ_j^\top$ transforms as $\bar S\mapsto R\bar SR^\top$ and the
noise is isotropic.
\item[(A2)] \emph{Permutation symmetry} of $P_\eta$ in the particle label.
\item[(A3)] \emph{Geometric ergodicity} with unique stationary law $\pi_\eta$ on $\cX^N$
(standard for uniformly elliptic noise on a compact state space).
\end{enumerate}
The diversity U-statistic is $\Delta=\frac{2}{N(N-1)}\sum_{i<j}h(V_i,V_j)$,
$h=\dG$, $\theta=\E_{\pi_2}[h(V_1,V_2)]$, with Hoeffding projections $h_1,h_2$ and
variance components $\zeta_1=\Var_{\pi_1}(h_1)$, $\zeta_2=\Var_{\pi_2}(h_2)$ as in the
problem statement.

% =====================================================================
\section{The invariance mechanism of the original conjecture fails}

\begin{proposition}[Exact marginal Haar; $h_1\equiv0$; $\zeta_1^*\equiv0$]
\label{prop:inv}
Under (A1)--(A3), for every $\eta\ge0$:
\begin{enumerate}
\item $\pi_\eta$ is $O(d)$-invariant (diagonal action) and exchangeable;
\item $\pi_{1}=\mu_H$, the Haar measure;
\item $m(x):=\E_{\pi_2}[h(V_1,V_2)\mid V_1=x]$ is constant in $x$, hence
$h_1\equiv0$, $\zeta_1=0$;
\item the ``effective variance'' of the problem statement,
$\zeta_1^*=\zeta_1+\sum_{j\ne1}\Cov_\pi(h_1(V_1),h_1(V_j))$, vanishes identically.
\end{enumerate}
\end{proposition}

\begin{proof}
(1) The generator $\mathcal L_\eta$ of $P_\eta$ commutes with the diagonal
$O(d)$-action by (A1): $T_R\mathcal L_\eta=\mathcal L_\eta T_R$. If $\mu$ is stationary
then $\mu T_R$ is stationary, so $\mu T_R=\mu$ by (A3). Exchangeability is identical
with the $S_N$-action.\\
(2) $\pi_1$ is an $O(d)$-invariant probability measure on the homogeneous space
$O(d)/(O(k)O(d-k))$; the Haar measure is the unique such measure (for any two points
$x,y$ choose $g$ with $gx=y$ and use $\mu(gA)=\mu(A)$ plus second countability to
identify the measure algebra), hence $\pi_1=\mu_H$.\\
(3) By diagonal invariance of $\pi_2$, for $g\in O(d)$:
\[
m(gx)=\E[h(gx,V')\mid V_1=gx]
=\E[h(gx,gV'')\mid V_1=x]=\E[h(x,V'')\mid V_1=x]=m(x),
\]
since $h(gx,gy)=h(x,y)$ and the conditional law of $V''=g^{-1}V'$ given $V_1=x$
equals that of $V'$ given $V_1=x$. Because $O(d)$ acts transitively on $\Gr(k,d)$,
$m$ is constant; $\E_{\pi_1} m=\theta$ gives $m\equiv\theta$, i.e.\ $h_1\equiv0$.\\
(4) Immediate from (3).
\end{proof}

\begin{corollary}
The mechanism asserted in the Invariance Dichotomy (``$\eta>0\Rightarrow\pi_1\ne\mu_H$'')
is false under (A1)--(A3); both sides of the variance formula in Problem 2
are identically zero and cannot explain any $\Theta(N^{-1})$ scaling.
The symmetry breaking produced by the consensus drift is \emph{spontaneous}:
conditioned on the (random) cluster location $F$, particles are distributed as
$\nu_F\ne\mu_H$, but $\pi_1=\int\nu_F\,\dd F=\mu_H$ by (2).
\end{corollary}

\begin{remark}[why $\eta>0$ nevertheless changes the pair structure]
While $\pi_1=\mu_H$ exactly, the two-marginal is
$\pi_2=\E_F[\nu_F\otimes\nu_F]\ne\mu_H^{\otimes2}$: two particles share their cluster
location. The dichotomy is a statement about $\pi_2$ and higher marginals, not about
$\pi_1$.
\end{remark}

% =====================================================================
\section{The dichotomy as a mean-field synchronization transition}

Let $L_N=\frac1N\sum_i\delta_{V_i}$ be the empirical measure and define, for a
probability measure $\mu$,
\[
\phi(\mu):=\iint h\,\dd\mu\,\dd\mu-\theta,
\qquad
m_\mu(x):=\int h(x,y)\,\dd\mu(y)-\theta .
\]
Note $\phi(\mu_H)=0$, $m_{\mu_H}\equiv0$, and $\Delta=\frac1N\sum_i g_i$ with the
leave-one-out means $g_i:=\frac1{N-1}\sum_{j\ne i}h(V_i,V_j)$.

\begin{lemma}[exact variance identity]
\label{lem:varid}
$\Var_\pi(\Delta)=\frac1{N^2}\Bigl[\sum_i\Var(g_i)+\sum_{i\ne j}\Cov(g_i,g_j)\Bigr]$.
\end{lemma}
\begin{proof}
$\Delta-\E\Delta=\frac1N\sum_i(g_i-\E g_i)$ and linearity.
\end{proof}

The dichotomy is a statement about which of the two sums dominates.

\begin{theorem}[disordered/chaotic regime; $\eta<\eta_c$]
\label{thm:dis}
Assume $\eta$ lies in the basin where $\mu_H$ is the stable mean-field equilibrium
(in particular $\eta=0$), and that the $m$-point marginals satisfy propagation of
chaos at rate $O(1/N)$ with the additional property that the directing-measure
fluctuations enter $\phi$ only at second order (which holds because
$h_1^{\mu_H}\equiv0$ kills the linear term). Then:
\begin{enumerate}
\item $\Var_\pi(\Delta)=\dfrac{2\,\zeta_2(\eta)}{N(N-1)}+\mathcal O(N^{-3})
=\Theta(N^{-2})$, where $\zeta_2(\eta)\to\Var_{\mu_H\otimes\mu_H}(h)=:\zeta_2(0)$;
\item $N\,(\Delta-\theta)\;\xto\;\displaystyle\sum_{\ell\ge1}\lambda_\ell\,
(\chi^2_{1,\ell}-1)$,
a (generally indefinite) weighted chi-square law, $\{\lambda_\ell\}$ the eigenvalues
of the centered kernel $h-\theta$ on $L^2(\mu_H)$; in particular
$\Var\bigl(N(\Delta-\theta)\bigr)\to2\sum_\ell\lambda_\ell^2=2\zeta_2(0)$ and
$\sqrt N(\Delta-\theta)\Rightarrow0$.
\end{enumerate}
\end{theorem}

\begin{proof}[Proof sketch (items (2) is classical; we indicate the $O(N^{-2})$ bound)]
(2) is the classical degenerate U-statistic CLT (Hoeffding; Serfling Thm.~5.5.1;
de Wet--Vitale for the infinite expansion): under $\mu_H$-i.i.d.\ sampling the
Hoeffding projection vanishes and $N(\Delta-\theta)$ is a V-statistic quadratic form
whose eigen-expansion yields the stated limit; tightness of the infinite series holds
because $\sum\lambda_\ell^2=\zeta_2(0)<\infty$. For the \emph{coupled} disordered
chain, write the law as a mixture $\pi_N\approx\int\nu^{\otimes N}\dd\Chi(\nu)$ with
directing measure $\Chi$ concentrated on an $O(N^{-1/2})$-ball around $\mu_H$
(mean-field fluctuation theorem). Conditional on the director, $\Delta-\theta=
\phi(\nu)+\frac2N\sum_i m_\nu(V_i)+W_N$ with $W_N$ the within-director degenerate
part, $\Var(W_N\mid\nu)=O(N^{-2})$. Since $m_\nu=m_{\mu_H}+(\partial m_\nu)=
O(\|\nu-\mu_H\|)=O(N^{-1/2})$ and
$\phi(\nu)=\phi(\mu_H)+\langle\phi'(\mu_H),\nu-\mu_H\rangle+O(\|\nu-\mu_H\|^2)
=0+2\langle h_1^{\mu_H},\nu-\mu_H\rangle+O(N^{-1}/N)=O(N^{-1})$
(the linear functional vanishes by Prop.~\ref{prop:inv}(3)),
$\Var\phi(\Chi)=O(N^{-2})$ and $\frac4N\E\Var(m_\nu)=O(N^{-2})$; the cross term is
$O(N^{-3/2}\cdot N^{-1/2})=O(N^{-2})$. Only the exact-$O(1)$ diagonal term
$2\zeta_2(\eta)/(N(N-1))$ survives, with $\zeta_2(\eta)$ a continuous perturbation of
the Haar value.
\end{proof}

\begin{theorem}[ordered/synchronized regime; $\eta>\eta_c$]
\label{thm:ord}
Assume the mean-field dynamics has a symmetry-broken stable equilibrium: the
stationary law is an equivariant mixture $\pi_N\approx\int\nu_F^{\otimes N}\dd F(F)$
with $F\sim\mu_H$ a random cluster location and cluster width
$w_\eta\asymp(D/\kappa\eta)^{1/2}$ ($\kappa$: curvature of the drift well). Then:
\begin{enumerate}
\item $\pi_1=\mu_H$, $h_1\equiv0$, $\zeta_1=0$ (Prop.~\ref{prop:inv}) --- the classical
non-degeneracy test gives a \emph{false negative};
\item nevertheless
$\Var_\pi(\Delta)=\dfrac{\sigma_\eta^2}{N}+\mathcal O(N^{-2})$ with
\[
\sigma_\eta^2 \;=\; 4\,\E_F\bigl[\Var_{\nu_F}(m_F)\bigr] \;+\; \gamma_\eta,
\qquad \gamma_\eta=O(N^{-1})\text{-channel},\quad
\E_F[\Var_{\nu_F}(m_F)]\asymp w_\eta^2\asymp \frac{D}{\kappa\eta};
\]
\item $\sqrt N\,(\Delta-\theta)\Rightarrow\mathcal N(0,\sigma_\eta^2)$;
\item there is no contradiction with the degeneracy: the $N^{-1/2}$-fluctuation is
carried by the \emph{conditional} first projection $m_F$, whose \emph{marginal}
average is killed by symmetry, while its conditional fluctuations are not.
\end{enumerate}
\end{theorem}

\begin{proof}[Proof sketch]
Condition on the director $F$ (de Finetti). Within the director the particle sample
is i.i.d.\ $\nu_F$, so
$\Delta-\theta=\phi(F)+\frac2N\sum_i m_F(V_i)+W_N(F)$ where now
$m_F\not\equiv0$: explicitly $m_F(x)=\E_{\nu_F}h(x,\cdot)-\theta$ is (by
equivariance) the pullback to $\cX$ of the geodesic distance to $F$, whose spread
over the cluster is $\Var_{\nu_F}(m_F)\asymp w_\eta^2$. The leading term:
$\Var\bigl(\frac2N\sum_i m_F(V_i)\mid F\bigr)=\frac4N\Var_{\nu_F}(m_F)$; averaging
over $F$ and adding the (subleading) $\Var\phi(F)=\gamma_\eta/N$ and the
$O(N^{-2})$ degenerate term gives (2). For (3), conditionally on $F$ the triangular
array $\{m_F(V_i)\}$ is i.i.d.\ with finite variance bounded away from zero, so the
martingale/Stein argument for exchangeable pairs (Chatterjee; or mixing CLT) gives
Gaussianity with variance $4\E_F\Var_{\nu_F}(m_F)$; the mixing hypothesis (A3) bounds
the cross-correlations $\Cov(g_i,g_j)=O(1/N)$, which contribute $O(N^{-1})$ to
$\Var(\Delta)$ and are asymptotically negligible at the $\sqrt N$-scale.
\end{proof}

\begin{theorem}[dichotomy = bifurcation]
\label{thm:dich}
Under (A1)--(A3) and the standard mean-field assumptions (uniform ellipticity,
attractive drift) there is a critical coupling $\eta_c\in(0,\infty)$
(Kuramoto-type linear instability of $\mu_H$; for $\Gr(1,2)$ the model is a mean-field
Kuramoto system on the circle) such that:
\[
\eta<\eta_c:\ \Var(\Delta)=\Theta(N^{-2}) \quad\text{vs.}\quad
\eta>\eta_c:\ \Var(\Delta)=\Theta(N^{-1}),
\]
and $N\,\Var(\Delta)$ is a susceptibility estimator that jumps at $\eta_c$ and decays
as $\sigma_\eta^2\asymp 4D/(\kappa\eta)$ in the ordered phase.
\end{theorem}

\begin{remark}[on $\Gr(1,2)$]
For $k=d=2$, $\Gr(1,2)$ is the projective line, the kernel is
$\phi(\Delta)=\arccos|\cos\Delta|$, and the dynamics is a mean-field Kuramoto model
with nematic (sign-symmetric) coupling. The Fourier coefficients are analytic:
\[
\lambda_\ell=-\frac{\sqrt2}{\pi\ell^2}\ (\ell\ \text{odd}),\qquad
\lambda_\ell=0\ (\ell\ \text{even}),\ \ s\text{-modes}=0,
\]
hence $\sum_\ell\lambda_\ell^2=\pi^2/48$ and
$\Var(N(\Delta-\theta))\to\pi^2/24$.
Our Monte Carlo gives $\sum\lambda_k^2=0.205617=\pi^2/48$ (machine agreement) and
empirical variances $0.4328$ ($N=300$) and $0.4048$ ($N=2000$) against the limit
$0.4112$ ($-4\%$/$+5\%$ finite-$N$ corrections, KS distance to the simulated weighted
chi-square $0.0091/0.0078$, below the $1.36/\sqrt{B}\approx0.0096$ noise floor).
\end{remark}

% =====================================================================
\section{Numerical verification (RTX 3060 campaign)}

Microscopic model (all $O(d)$-equivariant, permutation-symmetric):
frames $Q_i\in\mathbb R^{d\times k}$, per step of size $\dd t$,
\[
Q_i\leftarrow \mathrm{polar}\Bigl(Q_i+\dd t\,\eta\,2\bigl(\bar SQ_i-Q_i\,\mathrm{sym}(Q_i^\top\bar SQ_i)\bigr)
+\sqrt{2D\,\dd t}\,\xi_i\Bigr),
\]
Newton--Schulz polar retraction; $D=1$, $\dd t=\min(0.05,\;0.25/\eta)$;
$\Delta$ computed by exact pairwise principal angles (GPU batched SVD / closed form
for $k\le2$). Burn-in verified flat in $r$-traces.

\begin{center}
\begin{tabular}{lccc}
\hline
Regime & $N$-range & slope of $\log\Var$ vs $\log N$ & constant \\
\hline
$\eta=0$ (Haar, $\Gr(1,8)$) & 64--2048 & $-1.994$ & $N^2\Var\to0.098$--$0.105$\\
$\eta=5<\eta_c$ ($\Gr(1,8)$) & 64--2048 & $-2.013$ & $N^2\Var\to0.113$--$0.120$\\
$\eta=100>\eta_c$ ($\Gr(1,8)$) & 128--2048 & $-1.004$ & $N\Var\to6.65$--$6.79\times10^{-3}$\\
$\eta=100>\eta_c$ ($\Gr(2,8)$) & 128--512 & $-1.001$ & $N\Var\to7.03$--$7.11\times10^{-3}$\\
\hline
\end{tabular}
\end{center}

\textbf{Gaussianity, ordered phase} ($\eta=100$): pooled per-chain
$\sqrt N(\Delta-\theta)$: $|$skew$|\le0.13$, kurtosis $2.65$--$3.44$, QQ correlation
$\ge0.9988$ for all $N$; refutes a chi-square-type limit and confirms
Theorem~\ref{thm:ord}(3).

\textbf{Constant check:} measured $c_1$-proxy (within-swarm variance of the
leave-one-out means $g_i$) $=0.0016$; prediction $\sigma_\eta^2\approx4c_1=6.4\times
10^{-3}$ vs.\ measured $N\Var=6.65$--$6.79\times10^{-3}$: the conditional-projection
channel explains $\ge95\%$ of the variance.

\textbf{Transition:} at $N=512$, $r=\|\bar S-(k/d)I\|_F$ jumps
$0.063\to0.731$ between $\eta=24$ and $\eta=28$; $N\Var$ jumps
$1.23\times10^{-3}\to3.53\times10^{-2}$ (factor $\approx30$) and decays
$\propto\eta^{-1}$ ($\eta=100{:}\,6.6\times10^{-3}$, $200{:}\,3.2\times10^{-3}$),
matching $\sigma_\eta^2\asymp4D/(\kappa\eta)$.

\textbf{Haar limit law} ($\Gr(1,2)$): weighted chi-square confirmed
(Section 3, remark); strongly non-Gaussian (skew $-2.65$, kurtosis $13.8$), i.e.\
$\sqrt N(\Delta-\theta)\Rightarrow0$ there. High-statistics rerun ($B=40{,}000$,
kernel truncation $L=2000$): $N\in\{500,1500,4000\}$ gives variances
$0.407/0.415/0.429$ vs.\ limit $\pi^2/24=0.411$ and KS distances
$0.0060/0.0053/0.0068$, all below the $1.63/\sqrt B\in[0.008,0.012]$ noise floor.

\textbf{Decomposition check (campaign L3, $N=1024$):} at $\eta\in\{60,80,100\}$ the
conditional-projection channel $4c_1$ explains $97$--$99.6\%$ of $N\Var(\Delta)$
(e.g.\ $\eta=80$: $4c_1=8.255\times10^{-3}$ vs.\ $N\Var=8.29\times10^{-3}$), and
$\eta\,\sigma_\eta^2\approx0.64$--$0.70$ is constant, i.e.\
$\sigma_\eta^2\asymp D/(\kappa\eta)$ as predicted; the Gaussian limit persists
across the ordered phase (campaign L5: $\eta\in\{60,150\}$, $N=1024$:
$|$skew$|\approx0.14$, QQ--$r=0.998$).

\textbf{Coupled-disordered high stats (campaign L4):} at $\eta=5$, $N$ up to $4096$:
$N^2\Var=0.107$--$0.113$ (constant, slope $-2$), so Theorem~\ref{thm:dis} holds
far beyond the Haar point.

\textbf{Critical region (campaign L2):} at $\eta=28$, $N=1024$, the pooled ensemble
has disordered $\hat\theta=1.256$ but $N\Var=0.296$ (ten times the ordered-phase
level): individual chains tunnel between the disordered and ordered branches on the
sampling window --- the Ginzburg coexistence window, which narrows with $N$
($N=256$ is already ordered at $\eta=25$ while $N\ge512$ is not). $N\Var(\Delta)$ is
a sensitive inter-branch detector.

\textbf{Critical coupling scaling:} $\eta_c\propto D$ exactly
($\eta_c=13/25/50$ for $D=0.5/1/2$ at $d=8$) and $\eta_c$ grows approximately
linearly in $d$ ($12/25/72$ for $d=4/8/16$ at $D=1$; scaled by $d$: $3.0/3.1/4.5$
--- mild upward curvature at $d=16$, consistent with $\eta_c\sim D\,\kappa_{\ell^*}$
with $\kappa_2=2d$ plus higher-order gain corrections; finer grids and larger $N$
are queued).

% =====================================================================
\section{Corrections to the formal problem statement}

\begin{enumerate}
\item \textbf{Conjecture 1 (Invariance Dichotomy), part 2: \emph{invalid}} as stated.
$\pi_1=\mu_H$ for all $\eta$ (Prop.~\ref{prop:inv}); the observed dichotomy is real
but governed by the bifurcation (Thm.~\ref{thm:dich}), with a critical coupling
$\eta_c>0$, not by $\eta=0$ vs.\ $\eta>0$. Numerically: $\eta=5>0$ still gives slope
$-2.013$.
\item \textbf{Problem 2 (effective variance $\zeta_1^*$): vacuous.}
$\zeta_1^*\equiv0$. The correct leading coefficient in the ordered phase is
$\sigma_\eta^2=4\E_F[\Var_{\nu_F}(m_F)]+\gamma_\eta$, expressible through the
mean-field fluctuation field around the symmetry-broken equilibria $\nu_F$
(this vindicates the mean-field/Gaussian-field cycle-1 hypothesis).
\item \textbf{Problem 3 (CLT): holds in corrected form.}
$\sqrt N(\Delta-\theta)\Rightarrow\mathcal N(0,\sigma_\eta^2)$ for $\eta>\eta_c$
\emph{despite} $\zeta_1=0$; for $\eta<\eta_c$ the correct scaling is $N$ with
weighted-chi-square limit (confirmed; analytic constants on $\Gr(1,2)$).
\item \textbf{Practical corollary:} $N\Var(\Delta)$ is a synchronization
susceptibility estimator; the within-swarm variance of leave-one-out mean distances
$g_i$ is a plug-in estimator of $c_1$ and yields finite-sample confidence intervals
for $\sigma_\eta^2$ without bootstrap.
\end{enumerate}

\section{Follow-up cycle: analytic threshold and transition character}

The follow-up campaign (single-particle susceptibility, hysteresis/branch-persistence,
direct growth-rate measurement) resolved the bifurcation structure quantitatively.

\subsection{Analytic critical coupling}

The linear response of one particle to a static nematic field of amplitude $h$
(drift $2\eta h[(u\cdot v)u-(u\cdot v)^2v]$, measured, not assumed) gives
\[
\chi(\eta)\;=\;\frac{\dd a}{\dd h}\;\approx\;\frac{2\,(d-1)}{D\,d^2(d+2)}\;\eta,
\qquad a:=\E[(u\cdot v)^2]-\frac1d ,
\]
matching the Kubo-type prediction
$\langle\beta\cdot\nabla\psi_2\rangle/(D\kappa_2)$,
$\kappa_2=2d$, to \emph{four digits} at $d=8$
($\chi/\eta$: measured $0.0218$, theory $2\cdot7/(64\cdot10)=0.0219$).
The swarm supplies the field $h=c$ with $c=a/(1-1/d)$ the anisotropy coefficient of
$\E[vv^\top]-\I/d$, so the disordered branch loses stability at
\[
\boxed{\;\eta_c^{\mathrm{MF}}(d)\;=\;\frac{D\,d(d+2)}{2}\;}
\]
(independent derivations at $d=2$ via the exact von-Mises self-consistency
$I_1/I_0$ give the same value $4D$).
Measured gain crossings ($\chi=(d-1)/d$):
$d=2$: $\eta_c^{\mathrm{lin}}\approx4.5\pm0.5$ (theory $4$);
$d=4$: $\approx13$--$14$ (theory $12$);
$d=8$: $\approx42$ (theory $40$).

\subsection{Character of the transition: continuous at $d=2$, subcritical for $d\ge4$}

At $d=2$ the quartic drift term has no tangential component (pure Adler dynamics on
the folded circle), the bifurcation is the classical continuous one, and indeed the
ordered-branch persistence decays for every $\eta<\eta_c$ with no coexistence window.
For $d\ge4$ the branch-persistence and nucleation experiments reveal a genuine
\emph{coexistence window} $[\eta_A(d),\,\eta_c^{\mathrm{MF}}(d)]$:

\begin{center}
\begin{tabular}{lccc}
\hline
$d$ & ordered spinodal $\eta_A$ (measured) & $\eta_c^{\mathrm{MF}}$ (theory / measured) & character\\
\hline
$2$ & $=\eta_c$ (no window)        & $4$ / $\approx4.5$ & continuous\\
$4$ & $10$--$12$                   & $12$ / $\approx13$--$14$ & weakly first-order\\
$8$ & $19$--$20$                   & $40$ / $\approx42$ & strongly first-order\\
$16$& (below $72$, not resolved)   & $144$ / $\gtrsim110$ (noisy) & first-order (inferred)\\
\hline
\end{tabular}
\end{center}

Direct swarm growth-rate measurements ($N=2048$, $\log$-slope of early $r(t)$):
$\lambda\approx0.09/0.14/0.19$ at $\eta=25/28/30$ (noise-level wandering, no
ordering --- disordered branch still attractive) versus
$\lambda=1.3/2.6/4.4$ at $\eta=33/36/40$ with genuine ordering:
the finite-$N$ scans at $\eta\approx25$ ($d=8$) and $\eta\approx12$ ($d=4$)
sit \emph{inside} the coexistence window and are nucleation-controlled, with a
barrier growing with $N$ (explains: $N=256$ orders at $\eta=25$ while $N\ge512$
does not; ordered-$N=512$ nucleates at $\eta=30$ within $0.4$ time units while
disordered-$N=2048$ persists $15$ time units at the same $\eta$).

\subsection{Consequences for the diversity statistics}

\begin{enumerate}
\item The $\Theta(N^{-1})$ variance regime sets in as soon as \emph{any} branch of
the ordered phase is visited (metastability included): the dichotomy detector
$N\Var(\Delta)$ responds to branch occupation, not to the equilibrium phase
boundary --- finite-$N$ ``transitions'' measured by $\Delta$ are nucleation
diagnostics.
\item The bimodality of $\Delta$ inside the window (unimodal wide in disorder
$\theta\approx1.26$, unimodal low in order $\theta\approx0.30$--$0.61$, two-scale
mixtures in between) makes $\Delta$ a natural reaction coordinate for the
transition.
\item The analytic threshold $\eta_c^{\mathrm{MF}}(d)=D\,d(d+2)/2$ now
\emph{predicts} the $d$-dependence measured in the $\eta_c$-scan
($12/25/72$ for $d=4/8/16$ at $N=200$--$512$): each value lies inside its
coexistence window.
\end{enumerate}

\section{Open follow-ups}
\begin{enumerate}
\item Rigorous rate: replace the sketch in Thm.~\ref{thm:dis} by a full
Dirichlet-form/Talagrand argument; quantify $\gamma_\eta$ via the linearized
McKean--Vlasov fluctuation operator.
\item Critical exponents and nucleation-barrier scaling $N\Delta f(\eta)$ inside the
coexistence window (Eyring/Kramers fits to first-passage times).
\item Landau expansion of the mean-field free energy in powers of $c$ to prove the
subcriticality onset between $d=2$ and $d=4$ analytically (the sign of the cubic
term $\propto$ third cumulant of $(u\cdot v)^2$ under Haar).
\item Temporal FCLT for $\Delta_t$ (cycles 23/24) using the same directed-chain
machinery.
\end{enumerate}

\end{document}
